\begin{minipage}{0.65\linewidth}
    On considère le système oscillant (Solide-Ressort) où $m$ désigne la masse
    du solide et $K$ la raideur du ressort.

    On étudie le mouvement de $G$ dans un repère $(O,\ex)$ lié à la Terre
    supposé galiléen (Fig.~\ref{fig:solide_ressort}).

    On écarte $\mathrm{(S)}$ de sa position d'équilibre d'une distance $X_m$
    dans le
    sens opposé de $\ex$ puis on l'abandonne sans vitesse initiale à l'instant
    ($t_0=0$). Le solide $\mathrm{(S)}$ oscille sans frottement sur le plan
    horizontal.
\end{minipage}
\hfill
\begin{minipage}{0.34\linewidth}
    \begin{figure}[H]
        \centering
        \input{figures/solide_ressort}
        \caption{}
        \label{fig:solide_ressort}
    \end{figure}
\end{minipage}

L'équation horaire du mouvement de $G$ est~:
${x(t)=\num{5e-2}\cdot\cos(4\cdot t+\pi)}$ (\unit{\meter}).
\begin{donnees}
    \begin{itemize}
        \item $m=\qty{1.25}{\kg}$
    \end{itemize}
\end{donnees}
\begin{enumerate}
    \item Déterminer la valeur de la période propre $T_0$.
    \item Déduire la valeur de la raideur $K$.
    \item Déterminer la valeur de $\dot{x}_G$ vitesse de $G$ au passage par la
          position d'équilibre pour la première fois dans le sens positif.
\end{enumerate}
